\documentclass[10pt,a4paper]{amsart}
\usepackage[margin=19mm]{geometry}
\usepackage{amsmath,amssymb,mathtools}
\usepackage[scr=rsfs,cal=euler]{mathalfa}
\usepackage{xfrac}
\usepackage[unicode,hidelinks,bookmarks=false]{hyperref}
\newcommand{\ie}{i.e.\ }
\newcommand{\cf}{cf.\ }
\newcommand{\resp}{respectively\ }
\newcommand{\Subsection}{\subsection}
\providecommand{\nobreakdash}{\nobreak\hbox{-}}
\setlength{\emergencystretch}{2em}
\multlinegap0pt
\usepackage{amssymb}
\usepackage{xcolor}
\usepackage{enumerate}
\usepackage{tikz}
\newtheorem{lemma}{Lemma}
\newtheorem{proposition}{Proposition}
\title{An inverse problem for semilinear wave equations on metric tree graphs}
\author{Sergei Avdonin, Matti Lassas, Jinpeng Lu, Medet Nursultanov, Lauri Oksanen}
\date{Source: arXiv:2603.26170v1 (CC BY 4.0)}
\begin{document}
\maketitle

\noindent\emph{Source and licence.} An inverse problem for semilinear wave equations on metric tree graphs. Version arXiv:2603.26170v1 (CC BY 4.0), distributed by arXiv under the Creative Commons Attribution 4.0 International licence, http://creativecommons.org/licenses/by/4.0/, as recorded in the arXiv licence metadata for this exact version and shown on the abstract page https://arxiv.org/abs/2603.26170v1. Full licence notice: \texttt{licenses/arxiv-2603.26170-CC-BY-4.0.txt}.\par\smallskip

\noindent\textit{Editorial scope.} Counted unit: the article's lemma star-interfere (source0.tex lines 747--754). The article's complete original proof of it is delivered with it (source0.tex lines 756--860, one \texttt{proof} environment of 105 lines). No mathematics is written, repaired or re-derived: the statement and the proof are the article's own lines, in the article's own notation, and no retained line is substituted or re-flowed.  Retained uncounted as the source-local context the proof needs: lines 155--164; lines 172--174; lines 357--378; lines 358--364; lines 418--427; lines 538--563; lines 673--680; lines 692--697.  One declared adaptation: exactly 2 ancillary strings inside the retained spans are omitted (image-inclusion commands and, where the source repeats a label, the repeated label definitions) and no other line differs from the source; the figure environments, with their placement, captions and labels, are kept, and the package records the omitted strings in fidelity receipts bound to the affected bodies.  No retained line contains a citation, so the wrapper carries no bibliography and no bibliography entry.  The retained lines contain the article's own diagram code, which the wrapper typesets with the standard tikz packages; no image file is delivered and the package declares no asset dependency.  Editorial formatting only, in addition: an AMS \texttt{amsart} preamble that restates the article's own declarations and macros used by the retained lines, namely the environments \texttt{lemma}, \texttt{proposition} on separate counters, together with the standard packages those lines need.  The retained environments print in this document as Lemma 1, Proposition 1 in order of appearance, whereas the article prints the same items with the numbering of its own full document.  The complete native source of the article as distributed in the arXiv e-print for this exact version is bundled unchanged as \texttt{sources/197262/source0.tex}.
\begin{equation}\label{eq-wave-nonlinear}
    \begin{cases}
    (\partial_{t}^2 -\partial_x^2 ) u(x,t)+q(x) u(x,t)+a(x,t)u^3(x,t)=0, &x\in \Omega\setminus V, \ t \in (0,T),\\
    u(\cdot,t)|_{\Gamma_0}=f(t),\quad u(\gamma_0,t)=0 & t \in (0,T),\\
    %u(\gamma_0,t)=0, & \gamma_0\in \Gamma, \, t\in (0,T),\\
    u(x,t)|_{t \leq 0}=0, &  x\in \Omega, \\
%          u(x,0) = 0, \ \partial_t u(x,0) = 0, &x\in \Omega,\\
          u \textrm{ subject to the Kirchhoff-Neumann condition on } V\setminus\Gamma.
    \end{cases}
\end{equation}

\begin{equation}\label{DN-nonlinear}
\Lambda_T f = \partial u^f|_{\Gamma_0 \times (0,T)},
\end{equation}

This section is for the explicit formula, modulo $O(h)$, for the geometric optics solution of the (linear) wave equation on metric graphs with time-independent potential function $q$. We use this solution in the next section to find the coefficient $a(x,t).$
\begin{equation}\label{eq-gauss}
    \begin{cases}
    (\partial_{t}^2 -\partial_x^2 ) v(x,t)+q(x) v(x,t)=0, &x\in \Omega \setminus V, \  t > 0,\\    v(\cdot,t)|_{\Gamma}=f(t), & t> 0,\\
          v(x,t)|_{t\leq 0}=0, &x\in \Omega,\\
    v \textrm{ subject to the Kirchhoff-Neumann condition on } V\setminus\Gamma.
    \end{cases}
\end{equation}
The boundary source of our interest is of geometric optics type:
\begin{equation*}
\exp \big(i h^{-1} \theta(t) \big) \chi_b(t),
\end{equation*}
where $h>0$ is a small parameter, and $\chi_b$ is a smooth cutoff function near $b$ with an independent parameter $b>0$, say $\chi_b=1$ on $[b/2,3b/2]$ and vanishing outside of $(0,2b)$.
%The cutoff function $\chi_b$ is shifted a little so that the initial condition of the wave equation is satisfied.
To simplify computations, we choose 
\begin{equation} \label{theta}
\theta(t)=t,
\end{equation}
and the boundary source
\begin{eqnarray} \label{source}
f_h(t) := e^{ih^{-1}t} \chi_b(t).
\end{eqnarray}

\begin{equation}
    \begin{cases}
    (\partial_{t}^2 -\partial_x^2 ) v(x,t)+q(x) v(x,t)=0, &x\in \Omega \setminus V, \  t > 0,\\    v(\cdot,t)|_{\Gamma}=f(t), & t> 0,\\
          v(x,t)|_{t\leq 0}=0, &x\in \Omega,\\
    v \textrm{ subject to the Kirchhoff-Neumann condition on } V\setminus\Gamma.
    \end{cases}
\end{equation}

\begin{lemma} \label{CGO-Y}
On the metric graph of gluing $n$ intervals at one common point parametrized as above, the geometric optics solution of the equation \eqref{eq-gauss}, with the boundary source chosen as above, has the following form on the time interval $t\in [0,\min_{j} l_j]$:
\begin{equation}
    v_1^h(x,t)= e^{ih^{-1}(t+x)}\chi_b(t+x) - \frac{n-2}{n} e^{ih^{-1}(t-x)}\chi_b(t-x) +R_h(x,t), \quad x\in I_1, 
\end{equation}
\begin{equation}
    v_j^h(x,t)= \frac{2}{n} e^{ih^{-1}(t-x)}\chi_b(t-x) +R_h(x,t), \quad x\in I_j,\  j=2,\cdots, n,
\end{equation}
where $v_j^h$ is the restriction of the solution $v^h$ on $I_j$.
\end{lemma}

\begin{equation} \label{splitgroup}
\deg(v_i)\, g_i'(t) =F_i(t),
\end{equation}
 where for all $i=1, \dots, |V|-|\Gamma|$, $F_i(t)$ depends on the vector function $g$ with arguments delayed by some positive number. 
 The values of $g(t)$ on $\Gamma$ are given as boundary conditions in \eqref{eq-gauss}. Since $g(0)=0$, $g$ on $V \setminus \Gamma$ can be calculated in steps. 
 
%To compute $\psi_j(x),$ the first order terms in $h,$ on the edge $e_j$ we assume that our graph is tree, the control $f=f_h(t)$ is applied at one of the boundary vertices, say $v_1,$ and at all other  boundary vertices we have zero Dirichlet conditions. Let the edge $e_1(v_1,v_2)$ be incident to $v_1,$
%the edge $e_2(v_2,v_3)$ be incident to $v_2,$ and the edge $e_3(v_3,v_4)$ be incident to $v_3,$ and every $e_j(v_j,v_{j+1})$ is parametrized by $x \in (0,l_j)$ with $0$ at $v_j.$
%Then, as we demonstrated for a star graph,
 %$$\psi_1(x)=\int_0^x q_1(s)\,ds, \ \psi_2(x)=\int_0^x q_2(s)\,ds + \int_0^{l_1} q_1(s)\,ds, $$$$ \psi_3(x)=\int_0^x q_3(s)\,ds + \int_0^{l_2} q_2(s)\,ds +\int_0^{l_1} q_1(s)\,ds.  $$
 
% Since on a tree there is a unique path from $v_1$ to any other vertex, we can determine in such a way  $\psi_j(x)$ for every $e_i.$
 
 %{\bf S: I dropped here the factor $1/2$ of all integrals, it can be added later.
 %The result depends on our choice of the function $\theta$ and may have a slightly different form, but the approach works in any case.}


\section{Linearization method}
\label{sec-linearization}

Using the linearization method, the Dirichlet-to-Neumann map $\Lambda_T$ defined in \eqref{DN-nonlinear} determines
\begin{equation}
\partial_{\varepsilon_1} \partial_{\varepsilon_2} 
\partial_{\varepsilon_3}
\Big(\Lambda_T (\varepsilon_1 f_1+ \varepsilon_2 f_2+\varepsilon_3 f_3)\Big)  \Big|_{\varepsilon_1=\varepsilon_2=\varepsilon_3=0} = \partial w \big|_{\Gamma_0 \times (0,T)},
\end{equation}

\begin{lemma} \label{star-equal}
For a star graph of equal edge lengths, suppose that control and observation are available on one leaf $z$ (i.e., $\Gamma_0=\{z\}$).
%Dirichlet condition on all other leaves.
Let $e_z$ denote the edge incident to $z$ and $l$ be its length.
Assume $T>2l$.
Then $\Lambda_T$ uniquely determines $a$ on the edge $e_z$ in the domain
$$\Big\{(x,t)\in e_z\times [0,T]: 2l-d(x,z)\leq t\leq T-d(x,z)\Big\}.$$
\end{lemma}

\begin{proposition} \label{star-diff}
For any graph, suppose that control and observation are available on a vertex $z$ of degree 1 (i.e., $\Gamma_0=\{z\}$). {Let $e_z$ denote the edge incident to $z$ and $l$ be its length.
Assume $T>2l$.
Then $\Lambda_T$ uniquely determines $a$ on the edge $e_z$ in the same domain as in Lemma \ref{star-equal}.}
%$$\Big\{(x,t)\in e_0\times [0,T]: 2l-d(x,z_0)\leq t\leq T-d(x,z_0)\Big\}.$$
\end{proposition}

\begin{lemma} \label{star-interfere}
Consider a $3$-edge star graph where control and observation are available on two leaves $z_1,z_2$ (i.e., $\Gamma_0=\{z_1,z_2\}$), and let $z_0$ be the other leaf.
Denote by $e_i$ the edge incident to $z_i$ and by $l_i$ its length.
{Assume $l_1=l_2$ and $T>2(l_1+l_0)$.
Then $\Lambda_T$ uniquely determines $a$ on the domain
$$\big\{(x,t)\in e_j\times [0,T]: 2l_j-d(x,z_j)\leq t\leq T-d(x,z_j)  \big\},\quad j=1,2, \textrm{ and }$$
$$\big\{ (x,t)\in e_0\times [0,T]:  l_1+l_0+d(x,z_0) \leq t \leq T-l_1-l_0+d(x,z_0)\big\}.$$}
\end{lemma}

\begin{proof}
By Proposition \ref{star-diff}, $\Lambda_T$ uniquely determines $a$ on the edges $e_1,e_2$ in the domain stated in Lemma \ref{star-equal}. Hence we focus on the other edge $e_0$.
%where $z_0$ is the other boundary vertex where we do not have control.
%Denote by $e_i$ the edge incident to the boundary vertex $z_i$, $i=0,1,2$,
%and by $l_i$ the length of the edge $e_i$.
%We assume that $l_1=l_2$.
For convenience, we parametrize the edges $e_1$ and $e_0$ together as $[0,l_1+l_0]$, where $z_1$ corresponds to $x=0$; $z_0$ corresponds to $x=l_1+l_0$; $x=l_1$ corresponds to the internal vertex.
We focus on the behavior of waves on $e_0$, i.e., on $[l_1,l_1+l_0]$.


\begin{figure}
    \begin{minipage}{0.47\textwidth}
        
        \caption{The support of the wave $v_1$ (translated so that the wave starts from the origin).}
        \label{v1}
    \end{minipage}
    \begin{minipage}{0.47\textwidth}
        
        \caption{A two-step procedure for reconstructing $a$ on $[l_1,L]$. The blue lines represent the support of $v_2,v_3$.}
       \label{procedure}
    \end{minipage}
\end{figure}


Let $L=l_1+l_0$ so that $T>2L$ by assumption.
As before, we send geometric optics solutions $v_0,v_1,v_2,v_3$ from $z_1$, where $v_1$ is sent at time $t=t_0>0$; $v_2,v_3$ are sent at a delayed time $s>t_0$;
$v_0$ is sent at time $t=2L+t_0$.
What we need to figure out is where these four waves intersect.
{Recall that the requirement of $t_0>0$ is due to the initial condition of the equation \eqref{eq-wave-nonlinear}.
However, for the sake of explaining our construction, it is enough to demonstrate for $t_0=0$.}
We focus on the region over $e_0$, i.e., $$U_0:=[l_1,L]\times \mathbb [0,2L],$$ 
where the second coordinate is the time coordinate.

The support of the wave $v_1$ up to time $2L$ is illustrated in Figure \ref{v1}, where multiple reflections occur. Since we assume $l_1=l_2$, the reflected waves (grey lines) always come into $e_0$ in $2l_1$-increment.
In $U_0$,
$$\bigcap_{b>0}\textrm{supp}(v_1) \cap U_0 \subset \bigcup_{m\in \mathbb{N}} \{t-x=2m l_1\} \cup \{t+x=2L+2ml_1\}.$$
Similarly,
$$\bigcap_{b>0}\textrm{supp}(v_0) \cap U_0 \subset \bigcup_{m\in \mathbb{N}} \{t-x=-2m l_1\} \cup \{t+x=2L-2ml_1\}.$$
Recall that the waves are supported on a thin neighborhood of size $b$, where $b$ is the parameter in the cutoff function, see \eqref{source}.
Although these two supports above intersect frequently, one can see that the intersection is actually two lines union a finite number of points:
\begin{equation}\label{supp-01}
\bigcap_{b>0}\textrm{supp}(v_0) \cap \textrm{supp}(v_1) \cap U_0 =\{t=x\}\cup \{t+x=2L\} \cup \hat{A},
\end{equation}
where $\hat{A}$ is a finite set of points in $U_0$.
This gives much clarity about the intersections of four waves when we take the delayed waves $v_2,v_3$ into account as follows.
%However, the intersection on $[0,l_1]$ can be complicated. For example when $l_1$ divides $L$, the waves $v_0,v_1$ coincide on $[0,l_1]$, and the four waves can intersect at lines when delayed wave $v_2$ is sent at those specific delayed time. For our purpose, one can very well excludes these delayed time and recover them by continuity. In this lemma, one can include those specific delayed time, as the information required is included in the domain of recovery already obtained. However, in general construction, one has to be careful when the construction is done one step further.


\smallskip
First, we consider $v_2,v_3$ with delayed time parameter $s$ close to $2L$. 
If $2l_0=2L-2l_1<s<2L$ and $v_2$ satisfies
\begin{equation}\label{hatA}
\bigcap_{b>0}\textrm{supp}(v_2) \cap \hat{A} =\emptyset,
\end{equation}
%There exists a time $\hat{s}_0 \in (0,2L)$ such that if we send the waves $v_2,v_3$ with delayed time $$s\in (\hat{s}_0,2L), \quad \bigcap_{b>0}\textrm{supp}(v_2) \cap A_{01} =\emptyset,$$
there is no intersection region between $v_0,v_1,v_2,v_3$ in $U_0$ when $b$ is sufficiently small, see Figure \ref{procedure}.
In view of the finiteness of $\hat{A}$, the condition \eqref{hatA} excludes at most a finite number of $s$, since $\bigcap_{b>0} \textrm{supp}(v_2)$ is a finite collection of lines with slopes either $1$ or $-1$.

Next, we lower the time parameter $s$. There exists a time $\hat{s}_1 \in (0,2l_0)$ such that if we send the waves $v_2,v_3$ with delayed time $s\in (\hat{s}_1,2l_0)$ subject to \eqref{hatA},
%$$s\in (\hat{s}_1,2l_2), \quad \bigcap_{b>0}\textrm{supp}(v_2) \cap \hat{A} =\emptyset,$$
then there is only one intersection region between $v_0,v_1,v_2,v_3$ in $U_0$ when $b$ is sufficiently small, see Figure \ref{procedure}.
One can verify that $\hat{s}_1=2l_0-2l_1$ in the present situation, due to our assumption that $l_1=l_2$.
Since we assume to know the graph structure and edge lengths, we know the form of the waves by Lemma \ref{CGO-Y}, in particular, the coefficients depending on how the waves are reflected.
Then following the method in Section \ref{sec-linearization}, one can determine
\begin{equation} \label{lower-s-1}
C  a(L-\frac{s}{2},L+\frac{s}{2}) + \textrm{an integral over }e_1,e_2,
\end{equation}
where $C$ is some known constant. 
This integral over $e_1,e_2$ depends only on the pointwise values of $a$ at some points on $e_1,e_2$ where the four waves intersect.
%Consider if $L/l_1\in \mathbb{Z}$, the $v_1,v_0$ waves completely coincide on $e_1$. so the four waves should have multiple intersections on $e_1$.
On $e_1$, the waves $v_1,v_2$ can intersect only above (including) the line $x+t=2l_1$.
%it is enough to look at these two waves. $v_2$ is sent delayed than $v_1$.
On $e_2$, any intersection occurs only after time $t\geq l_1+d(x,v_S)=l_1+l_2-d(x,z_2)$, $x\in e_2$, where $v_S$ is the internal vertex of the star graph,
because this is the shortest travel time needed for the waves reaching $x\in e_2$ from $z_1$.
%The actual intersection time on $e_2$ is strictly larger than the shortest travel time (except endpoints). The following two situations may happen.
%(1) The intersection only occurs after reflection of $v_1$ on $z_2$. Then the time is at least $l_1+l_2$.
%(2) $v_2$ coincides on $e_2$ with the reflected wave of $v_1$ after hitting $z_0$. In this case, the intersection time is at least $l_1+2l_0+d(x,v_S)$.
On either $e_j$, $j=1,2$, the intersections with the wave $v_0$ must satisfy $t\leq 2L-d(x,z_1)$, $x\in e_j$.
%this is the shortest time to reach $z_1$ at time $2L$.
Recall that $T>2L$ and $d(x,z_1)\geq l_1=l_2\geq d(x,z_2)$ for $x\in e_2$.
Hence, the integral over $e_1,e_2$ in \eqref{lower-s-1} can be computed, since
$a$ is already recovered on $e_1,e_2$ in the domain $\{(x,t)\in e_j\times [0,T]: 2l_j\leq t+d(x,z_j)\leq T\}$ for $j=1,2$ due to Proposition \ref{star-diff}.
%Note that we have used the same parametrization on $e_2$ as $e_1$, namely $z_2$ corresponding to $x=0$.
%The parametrizations on $e_1,e_2$ are exactly the same.
This procedure determines $a$ on 
$$\big\{x+t=2L, \; x\in [l_1, L-\hat{s}_1/2]=[l_1,2l_1]\big\}$$ 
except for a finite number of points excluded by the condition \eqref{hatA}, which determines $a$ on $\{x+t=2L,\, x\in [l_1, 2l_1]\}$ by continuity. 

Note that if $l_1\geq l_0$ the procedure stops after the previous step. Otherwise,
we keep lowering $s$ up to some time $\hat{s}_2$ while avoiding the finite set $\hat{A}$, more precisely $\hat{s}_2=2l_0-4l_1$, and we have two intersection regions in $U_0$ for sufficiently small $b$. Thus one determines the sum
$$C  a(L-\frac{s}{2},L+\frac{s}{2}) + C'  a(x',2L-x') + \textrm{an integral over }e_1,e_2,$$
where $x'=L-s/2-l_1\in [l_1,2l_1]$ in our present situation. Since we have already solved $a$ for $x\in [l_1,2l_1]$, this determines $a$ on the line $\{x+t=2L\}$ for $x\in [2l_1,3l_1]$.
{Repeating this procedure determines $a$ on the line $\{x+t=2L,\, x\in [l_1,L]\}$ in finite steps.}
%Combining with Proposition \ref{star-diff}, we can determine $a$ on $\big\{ x+t=2L,\; x\in [0,L]  \big\}.$

Now instead of sending the wave $v_1$ from $t_0=0$, we vary $t_0$ in $(0,T-2L)$ and repeat the same construction above on the line $\{x+t=2L+t_0\}$. 
Proposition \ref{star-diff} is valid since $2l_1 < 2L+t_0< T$.
Then the same arguments above determine $a$ on 
$\big\{x+t=2L+t_0\in (2L,T), \, x\in [l_1,L] \big\}$, which determines $a$ on $\big\{2L \leq x+t \leq T, \, x\in [l_1,L] \big\}$ by continuity.
The latter domain can be translated to the following unparametrized form:
$$\big\{(x,t)\in e_0\times [0,T]: 
L+d(x,z_0) \leq t \leq T-L+d(x,z_0) \big\},$$
as stated in the lemma.
%Identifying $x$ with $L-d(x,z_0)$.
\end{proof}

\end{document}
